\section{Zeta and Mordell--Tornheim traces}
\label{sec:traces-determinants}

Define the coprime double series
\begin{equation}
 P(s)=\sum_{\substack{a,b\geq1\\(a,b)=1}}
 a^{-s}b^{-s}(a+b)^{-2s}.
 \label{eq:primitive-P}
\end{equation}
The majorant in \eqref{eq:hs-majorant} shows that it converges absolutely
for $\sigma>1/2$.  We first compute traces only where operator ideals make
them legal; analytic continuation plays no role.

\begin{theorem}[The two legal traces]
\label{thm:two-traces}
For $\sigma>1$,
\begin{equation}
 \Tr(E_s)=2^{-s}\zeta(s).
 \label{eq:first-trace}
\end{equation}
For $\sigma>1/2$,
\begin{equation}
 \Tr(E_s^2)=\zeta(2s)P(s)
 =\frac{\zeta(2s)}{\zeta(4s)}\zMT(s,s;2s).
 \label{eq:second-trace}
\end{equation}
Both statements are identities of absolutely convergent series in the
displayed domains.
\end{theorem}

\begin{proof}
If $\sigma>1$, \cref{prop:trace-wall} makes $E_s$ trace class, so its trace
is the absolutely convergent standard-basis diagonal sum.  The loops are
exactly $m=2t$, and therefore
\[
 \Tr(E_s)=\sum_{t\geq1}(2t)^{-s}=2^{-s}\zeta(s).
\]

Now suppose $\sigma>1/2$.  By \cref{prop:hs-wall}, $E_s$ is
Hilbert--Schmidt, hence $E_s^2$ is trace class.  Moreover,
\[
 \sum_{m,n}|e_s(m,n)e_s(n,m)|
 =\sum_{m,n}|e_s(m,n)|^2<\infty,
\]
so the trace may be evaluated term by term:
\begin{equation}
 \Tr(E_s^2)
 =\sum_{\substack{m,n\geq1\\m+n\mid mn}}(mn)^{-s}.
 \label{eq:ordered-edge-trace}
\end{equation}
This is an ordered-edge sum.  Applying
\cref{prop:coprime-coordinate} and using
$mn=t^2ab(a+b)^2$ gives
\[
 \Tr(E_s^2)
 =(\sum_{t\geq1}t^{-2s})
  \sum_{\substack{a,b\geq1\\(a,b)=1}}
    a^{-s}b^{-s}(a+b)^{-2s}
 =\zeta(2s)P(s).
\]
No factor of two is inserted: $(a,b)$ is already ordered.

Finally, decompose an arbitrary ordered pair $(u,v)$ uniquely as
$u=ga$, $v=gb$ with $(a,b)=1$.  Absolute convergence permits regrouping,
and \eqref{eq:mt-definition} becomes
\begin{align*}
 \zMT(s,s;2s)
 &=\sum_{g\geq1}g^{-4s}
   \sum_{\substack{a,b\geq1\\(a,b)=1}}
     a^{-s}b^{-s}(a+b)^{-2s}\\
 &=\zeta(4s)P(s).
\end{align*}
Substitution proves the second equality in \eqref{eq:second-trace}.
\end{proof}

The adjective ``coprime'' in $P(s)$ refers only to $(a,b)=1$.  The
arbitrary-pair sum $\zMT(s,s;2s)$ and the elementary gcd extraction remain
classical arithmetic objects; the graph-specific statement is that the
same operator whose even diagonal gives \eqref{eq:first-trace} has
\eqref{eq:second-trace} as its two-step trace.

\subsection{Determinant consequences}

The ideal walls determine exactly which determinant is available.  We use
the standard convention
\[
 \det_2(I-zA)=\det((I-zA)e^{zA})
 \qquad(A\in\Schatten{2}).
\]

\begin{corollary}[Legal determinants and local logarithms]
\label{cor:determinants}
If $\sigma>1/2$, the function
$D_2(z;s)=\det_2(I-zE_s)$ is entire in $z$.  On the disk
$|z|\,\|E_s\|<1$, with the logarithm normalized by
$\log D_2(0;s)=0$,
\begin{equation}
 \log D_2(z;s)
 =-\sum_{r\geq2}\frac{z^r}{r}\Tr(E_s^r),
 \label{eq:det2-local-log}
\end{equation}
and in particular
\begin{equation}
 [z^2]\log D_2(z;s)
 =-\frac12\frac{\zeta(2s)}{\zeta(4s)}
   \zMT(s,s;2s).
 \label{eq:det2-quadratic}
\end{equation}
If $\sigma>1$, the ordinary determinant
$D_1(z;s)=\det(I-zE_s)$ is entire in $z$ and
\begin{align}
 \log D_1(z;s)&=-\sum_{r\geq1}\frac{z^r}{r}\Tr(E_s^r),
 \label{eq:det1-local-log}\\
 D_2(z;s)&=D_1(z;s)\exp\!(z\Tr(E_s))
 \label{eq:det-overlap}
\end{align}
on the same sufficient local disk for the logarithm.
\end{corollary}

\begin{proof}
These are the standard Hilbert--Carleman and Fredholm determinant
identities for Schatten operators \citep{Simon2005}.  In the first domain,
$E_s^r$ is trace class for every $r\geq2$ because $E_s^2$ is trace class
and $E_s$ is bounded.  Substituting \eqref{eq:second-trace} in the
$r=2$ term gives \eqref{eq:det2-quadratic}.  In the trace-class overlap,
the definition of $\det_2$ gives \eqref{eq:det-overlap}, and
\eqref{eq:first-trace} supplies the exponent.
\end{proof}

Only the determinant functions are asserted globally in $z$.  The
logarithmic series and the chosen logarithm are local.  We do not use the
right side of \eqref{eq:second-trace} to continue a determinant past an
ideal wall, and we do not claim a completed divisor or functional equation.
